\section{Construction, Security Setting, and Result}
\label{sec:security-model}

This section defines the construction and states the result proved later,
explaining formalization-specific conventions as they arise.
All distributions are discrete.
Games return a Boolean, and
\(\WinProb[G]=\Pr[G=\mathsf{true}]\).
We write
\(\TVD(P,Q)=\tfrac12\sum_x|P(x)-Q(x)|\) for statistical distance and
\(\KL{P}{Q}\) for KL divergence when the defining series is summable and
\(P\) is absolutely continuous with respect to \(Q\).

\subsection{Approximate FHE and its flooded transform}

CKKS introduced homomorphic arithmetic with approximate plaintext
semantics~\cite{ckks}.
LMSS subsequently gave a generic approximate-correctness formulation in which
ciphertexts carry public error estimates~\cite{lmss}.
Our formal scheme interface is a typed, executable presentation of that
setting:
\[
 S=(\mathsf{KeyGen},\mathsf{Enc},\mathsf{Eval}_1,
       \mathsf{Eval}_2,\mathsf{Dec}).
\]
For bookkeeping, the formal model represents a ciphertext as either
\(\mathsf{None}\), for an unsupported evaluation, or
\(c=\mathsf{Some}(\bar c,e)\), pairing a ciphertext with its public error bound
\(e\in\bbN\).
This option wrapper belongs to the formal API, not to the cryptographic
construction.
Unary and binary evaluation gates suffice to describe arbitrary circuits while
keeping the package interface small.

% IEEEtran can place a double-column float only after it has been encountered.
% Keep this declaration early enough for the table to appear at the top of the
% next page; the FloatBarrier at the end of this section then causes no gap.
\begin{table*}[t]
\centering
\caption{Formal theorem boundary: inputs and obligations of the checked
security theorem.
The first
five rows are inputs to the Rocq theorem; the last row explains what its
semantic model does not express.}
\label{tab:theorem-boundary}
\begin{tabular}{@{}p{.18\textwidth}p{.34\textwidth}p{.39\textwidth}@{}}
\toprule
Input & What Rocq assumes & Consequence used by the proof \\
\midrule
Scheme \(S\) & Typed, distribution-valued algorithms; key generation returns
with probability one & Defines the IND-CPA game and noise-flooded scheme \\
Noise coordinates & Dimension \(n\), chart maps, and the three laws in
\Cref{eq:charts} & Transports product discrete-Gaussian noise from
\(\bbZ^n\) to messages and measures the displacement between its centers \\
Probability-one correctness & Deterministic decryption; every output having
nonzero probability satisfies \Cref{eq:approx-correctness} & Every reachable
table ciphertext is within its recorded error bound \\
IND-CPA security & For every adversary \(B\) with the required oracle interface,
\(\WinProb[\mathsf{IND\text{-}CPA}_S(B)]\leq\beta_{\mathsf{CPA}}(B)\) &
Supplies the bound for the final game \\
Flooding parameter & \(\gamma>0\) & Makes division and the Gaussian loss
well defined \\
Adversary model & A well-typed SSProve oracle program; running time is not part
of the semantics & Separate inspection that the explicit reduction preserves
probabilistic polynomial time~\cite{ssprove} \\
\bottomrule
\end{tabular}
\end{table*}

Although the interface permits randomized decryption, the theorem specializes
it to a deterministic value \(\mathsf{dec}_0(sk,c)\) and assumes
probability-one approximate correctness.
For every encryption or evaluation output
\(c=\mathsf{Some}(\bar c,e)\) having nonzero probability and representing
\(m\),
\begin{equation}
\label{eq:approx-correctness}
 d(\mathsf{dec}_0(sk,c),m)\leq e.
\end{equation}
The artifact calls this \emph{support-level} correctness; in cryptographic
terms, the bound holds with probability one.
Nonzero correctness error would require the additional up-to-bad hop discussed
in \Cref{sec:scope}.

LMSS's noise-flooding defense postprocesses approximate decryption with fresh
Gaussian noise~\cite{lmss}.
On a concrete plaintext ring, adding such noise to a plaintext is an ordinary
algebraic operation.
For an abstract message type \(M\), however, a product discrete Gaussian is
supported on \(\bbZ^n\), not on \(M\); it cannot be used as message-valued
noise without first supplying integer coordinates.

Our formalization therefore packages the metric and the coordinates needed
for noise flooding as one interface.
The message space carries a distance \(d:M\times M\to\bbN\), and for each
center \(a\in M\) it has maps
\[
 I_a:M\rightarrow\bbZ^n,
 \qquad J_a:\bbZ^n\rightarrow M
\]
satisfying
\begin{equation}
\label{eq:charts}
 I_a(a)=0,
 \quad d(a,b)=\norm{I_a(b)}_\infty,
 \quad J_b(v)=J_a(v+I_a(b)).
\end{equation}
Here \(I_a(b)\) is the integer displacement from \(a\) to \(b\), while
\(J_a(v)\) interprets the integer displacement \(v\) as a message based at
\(a\).
Thus \(J_a\) transports a distribution on \(\bbZ^n\) to a distribution on
messages.
The three laws say that the coordinates are centered, that their
\(\ell_\infty\) norm agrees with message distance, and that changing the base
point translates the integer coordinates.
In particular, they let the proof compare flooding around two messages as
centered and shifted Gaussians in one common copy of \(\bbZ^n\).

For \(\gamma>0\), noise flooding around \(x\in M\) is consequently the
following single construction:
\begin{equation}
\label{eq:flood}
\begin{aligned}
 \mathsf{Flood}_e(x):\quad
 \sigma_e&=\max\{1,e\}\gamma,\\
 \eta&\leftarrow D_{\bbZ^n,0,\sigma_e^2},\\
 \mathsf{return}\;&J_x(\eta).
\end{aligned}
\end{equation}
Here \(D_{\bbZ^n,0,\sigma_e^2}\) is the product of \(n\) independent
integer discrete Gaussians.
The sampling step lives in \(\bbZ^n\), and \(J_x\) makes its result a message;
an arbitrary metric space without such maps would not support
\Cref{eq:flood}.
This proof-specific interface is easier to audit than committing the generic
theorem to a concrete CKKS polynomial quotient ring.
The noise-flooded transform \(\mathsf{NF}_{\gamma}[S]\) keeps key generation,
encryption, and evaluation unchanged and decrypts a valid tagged ciphertext by
\[
 a\leftarrow\mathsf{Dec}(sk,c);
 \quad \mathsf{Flood}_e(a).
\]
On input \(\mathsf{None}\), decryption fails; SSProve represents this by a
subdistribution of total weight zero.

We use the standard multi-query left-or-right IND-CPA experiment and the
IND-CPAD experiment of Li and Micciancio~\cite{lm}; their full definitions
appear in \Cref{alg:ind-cpa-game,alg:ind-cpad-game,alg:ind-cpad-oracles}.
For the discussion below, the essential difference is that IND-CPAD adds
evaluation and restricted decryption: a decryption query is answered only
when the two plaintexts recorded for that ciphertext agree.
We write the assumed IND-CPA bound as
\[
 \WinProb[\mathsf{IND\text{-}CPA}_{S}(\cB)]
 \leq \beta_{\mathsf{CPA}}(\cB).
\]

\subsection{The one-query decryption replacement}
\label{sec:simulation-target}

The reduction compares two ways to answer a decryption query when the two
recorded plaintexts agree.
Fix a row that can occur with nonzero probability,
\[
 T_i=(m,m,c),
 \qquad c=\mathsf{Some}(\bar c,e).
\]
Under deterministic decryption, the real and simulated answers are
\begin{align}
\label{eq:real-sim-decrypt}
 \mathsf{ODec}_{\mathsf{real}}(i):\quad&
   y\leftarrow\mathsf{Flood}_e(\mathsf{dec}_0(sk,c));
   \ \mathsf{return}\ \mathsf{Some}(y),\notag\\
 \mathsf{ODec}_{\mathsf{sim}}(i):\quad&
   y\leftarrow\mathsf{Flood}_e(m);
   \ \mathsf{return}\ \mathsf{Some}(y).
\end{align}
Both versions perform the same bounds checks and counter updates, refuse the
same queries, and abort in the same circumstances.
The simulated answer needs neither \(sk\) nor \(b\): the common plaintext
\(m\) is already recorded in the table.

The probability-one correctness assumption in \Cref{eq:approx-correctness}
gives
\(d(\mathsf{dec}_0(sk,c),m)\leq e\).
By the change-of-basepoint law in \Cref{eq:charts}, the simulated answer can
be expressed in the real answer's chart by shifting the integer Gaussian by
\(I_{\mathsf{dec}_0(sk,c)}(m)\).
The distance law bounds the \(\ell_\infty\) norm of this shift by \(e\), so
each of its \(n\) coordinates has magnitude at most \(e\).
The equal-variance discrete-Gaussian identity
\begin{equation}
\label{eq:dg-kl}
 \KL{\DG{\mu}{\sigma^2}}{\DG{\nu}{\sigma^2}}
   =\frac{(\nu-\mu)^2}{2\sigma^2}.
\end{equation}
After fixing the adversary's prior view, \(m\), \(c\), and \(e\) are fixed.
The identity bounds the KL divergence \(\epsilon_i\) between the two answer
distributions for this call by
\begin{equation}
\label{eq:epsilon-nf}
 \epsilon_i
 \leq
 \frac{ne^2}{2\max\{1,e\}^2\gamma^2}
 \leq
 \epsilon_{\mathsf{nf}},
 \qquad
\epsilon_{\mathsf{nf}}=\frac{n}{2\gamma^2}.
\end{equation}

This is a one-call statement conditioned on an arbitrary prior interaction
history.
The central problem is to preserve it across the adversary's \(q\) adaptive
queries without converting to statistical distance after each call.

The underlying noise-flooding reduction is due to LMSS~\cite{lmss}.
Our contribution is to express its assumptions as the explicit interface in
\Cref{tab:theorem-boundary} and check the resulting adaptive reduction in Rocq.

\begin{theorem}[Checked noise-flooding reduction]
\label{thm:main}
Let \(S\), its noise-coordinate interface, probability-one correctness proof,
IND-CPA security interface, and \(\gamma>0\) satisfy
\Cref{tab:theorem-boundary}.
For every IND-CPAD adversary \(\cA\) represented as a well-typed SSProve oracle
program, and every \(q\in\bbN\), the formalization constructs an IND-CPA adversary
\(\cB_{\cA,q}\) with the required oracle interface and proves
\begin{align*}
 &\WinProb[
   \mathsf{IND\text{-}CPAD}_{\mathsf{NF}_{\gamma}[S],q}(\cA)]\\
 &\quad\leq \beta_{\mathsf{CPA}}(\cB_{\cA,q})
       +\sqrt{q\epsilon_{\mathsf{nf}}/2}\\
 &\quad\leq \beta_{\mathsf{CPA}}(\cB_{\cA,q})
       +\frac{\sqrt{qn}}{2\gamma}.
\end{align*}
\end{theorem}

The first inequality keeps the per-query bound
\(\epsilon_{\mathsf{nf}}\) visible; substituting
\(\epsilon_{\mathsf{nf}}=n/(2\gamma^2)\) from \Cref{eq:epsilon-nf} gives the
second.
Thus IND-CPA security is an assumption on \(S\), while IND-CPAD security of its
noise-flooded version is the conclusion.

\FloatBarrier
